\begin{corollary}[黄金分支 Fibonacci 审计常数的双相算术平均恰为 \texorpdfstring{$8/15$}{8/15}]\label{cor:xi-golden-w1-two-phase-arithmetic-mean-eight-fifteenths}
\leanverified{paper\_xi\_golden\_w1\_two\_phase\_arithmetic\_mean\_eight\_fifteenths}
沿用定理 \ref{thm:xi-golden-w1-true-two-phase-limit} 的记号。则黄金分支在 Fibonacci 分母长度上的两相极限常数之算术平均满足
\[
\boxed{\;
\frac12\lim_{\substack{n\to\infty\\ n\ \mathrm{even}}}
F_n\,\mathsf{T}_{F_n}(\varphi^{-1})
+
\frac12\lim_{\substack{n\to\infty\\ n\ \mathrm{odd}}}
F_n\,\mathsf{T}_{F_n}(\varphi^{-1})
=
\frac{8}{15}.\;}
\]
等价地，
\[
\boxed{\;
\frac12\left(\limsup_{n\to\infty}F_n\,\mathsf{T}_{F_n}(\varphi^{-1})
+
\liminf_{n\to\infty}F_n\,\mathsf{T}_{F_n}(\varphi^{-1})\right)
=
\frac{8}{15}.\;}
\]
\end{corollary}
\begin{proof}
由定理 \ref{thm:xi-golden-w1-true-two-phase-limit}，
\[
\lim_{\substack{n\to\infty\\ n\ \mathrm{even}}}
F_n\,\mathsf{T}_{F_n}(\varphi^{-1})
=
\frac12+\frac{1}{2\sqrt5},
\]
\[
\lim_{\substack{n\to\infty\\ n\ \mathrm{odd}}}
F_n\,\mathsf{T}_{F_n}(\varphi^{-1})
=
\frac12-\frac{1}{2\sqrt5}+\frac1{15}.
\]
两式取算术平均，立即得到
\[
\frac12\left(\frac12+\frac{1}{2\sqrt5}\right)
+
\frac12\left(\frac12-\frac{1}{2\sqrt5}+\frac1{15}\right)
=
\frac{8}{15}.
\]
又因该定理已给出两子列分别对应于全序列的 \(\limsup\) 与 \(\liminf\)，故第二个盒装公式只是同一恒等式的重述。证毕。
\end{proof}

\begin{corollary}[黄金分支 Fibonacci 传输常数的双相振幅与中心值]\label{cor:xi-golden-w1-two-phase-amplitude-center}
\leanverified{paper\_xi\_golden\_w1\_two\_phase\_amplitude\_center}
沿用定理 \ref{thm:xi-golden-w1-true-two-phase-limit} 的记号。则数列
\[
F_n\,\mathsf{T}_{F_n}(\varphi^{-1})
\]
不收敛，并且恰有两个聚点
\[
\frac12+\frac{1}{2\sqrt5},
\qquad
\frac12-\frac{1}{2\sqrt5}+\frac1{15}.
\]
特别地，
\[
\boxed{\;
\limsup_{n\to\infty}F_n\,\mathsf{T}_{F_n}(\varphi^{-1})
=
\frac12+\frac{1}{2\sqrt5},\;}
\]
\[
\boxed{\;
\liminf_{n\to\infty}F_n\,\mathsf{T}_{F_n}(\varphi^{-1})
=
\frac12-\frac{1}{2\sqrt5}+\frac1{15}.\;}
\]
若记双相振幅
\[
\Omega_\varphi
:=
\limsup_{n\to\infty}F_n\,\mathsf{T}_{F_n}(\varphi^{-1})
-
\liminf_{n\to\infty}F_n\,\mathsf{T}_{F_n}(\varphi^{-1}),
\]
则
\[
\boxed{\;
\Omega_\varphi
=
\frac1{\sqrt5}-\frac1{15}.\;}
\]
同时，两聚点的算术中心满足
\[
\boxed{\;
\frac12\left(\limsup_{n\to\infty}F_n\,\mathsf{T}_{F_n}(\varphi^{-1})
+
\liminf_{n\to\infty}F_n\,\mathsf{T}_{F_n}(\varphi^{-1})\right)
=
\frac{8}{15}.\;}
\]
\end{corollary}
\begin{proof}
由定理 \ref{thm:xi-golden-w1-true-two-phase-limit}，偶指标与奇指标 Fibonacci 子序列分别收敛到
\[
\frac12+\frac{1}{2\sqrt5},
\qquad
\frac12-\frac{1}{2\sqrt5}+\frac1{15},
\]
且二者不同，故全序列不收敛而恰有这两个聚点；相应的 \(\limsup\) 与 \(\liminf\) 公式随即成立。

两式相减得到
\[
\Omega_\varphi
=
\left(\frac12+\frac{1}{2\sqrt5}\right)
-
\left(\frac12-\frac{1}{2\sqrt5}+\frac1{15}\right)
=
\frac1{\sqrt5}-\frac1{15}.
\]
而两式的算术平均恰为推论
\ref{cor:xi-golden-w1-two-phase-arithmetic-mean-eight-fifteenths} 的内容，因此中心值为 \(8/15\)。证毕。
\end{proof}

\begin{theorem}[黄金分支在金属平均族中唯一最大化单位符号预算下的分母熵率]\label{thm:xi-golden-metallic-symbol-budget-entropy-optimality}
\leanverified{paper\_xi\_golden\_metallic\_symbol\_budget\_entropy\_optimality}
沿用定义 \ref{def:const-type} 与定理 \ref{thm:metallic-gap} 的记号。对每个整数 \(A\ge 1\)，令
\[
\alpha_A=[0;A,A,A,\dots],
\]
并记其第 \(n\) 个收敛分母为 \(q_n(A)\)。定义单位符号预算下的分母熵率
\[
\eta(A):=
\lim_{n\to\infty}\frac{1}{An}\log q_n(A).
\]
则极限存在，且有显式公式
\[
\boxed{\;
\eta(A)=\frac{\log\lambda_A}{A}=\kappa(A)^{-1}.\;}
\]
因此函数 \(A\mapsto \eta(A)\) 在 \(A\ge 1\) 上严格递减，并在黄金分支 \(A=1\) 处取得唯一最大值
\[
\boxed{\;
\eta_{\max}=\eta(1)=\log\varphi.\;}
\]
换言之，在常数型金属平均族中，黄金分支以最少的单位符号预算实现最大的分母指数增长率；任意 \(A\ge 2\) 都严格逊于 \(A=1\)。
\end{theorem}
\begin{proof}
常数型连分数的分母满足标准递推
\[
q_{n+1}(A)=Aq_n(A)+q_{n-1}(A),
\]
其特征多项式为
\[
X^2-AX-1.
\]
由定义 \ref{def:const-type}，该多项式的 Perron 根正是
\[
\lambda_A=\frac{A+\sqrt{A^2+4}}{2}.
\]
因此存在常数 \(c_A\neq 0\) 使
\[
q_n(A)=c_A\lambda_A^n+O(\lambda_A^{-n}),
\]
从而
\[
\lim_{n\to\infty}\frac{1}{n}\log q_n(A)=\log\lambda_A.
\]
两边再除以 \(A\)，便得到
\[
\eta(A)=\frac{\log\lambda_A}{A}.
\]
而定理 \ref{thm:metallic-gap} 已知
\[
\kappa(A):=\frac{A}{\log\lambda_A}
\]
在 \(A\ge 1\) 上严格递增，故
\[
\eta(A)=\kappa(A)^{-1}
\]
必严格递减，并在最小允许值 \(A=1\) 处取得唯一最大值。最后，由
\[
\lambda_1=\varphi
\]
即得
\[
\eta(1)=\log\varphi.
\]
证毕。
\end{proof}

\begin{corollary}[金属平均族的逐位指数效率由黄金端点唯一极大化]\label{cor:xi-golden-metallic-perdigit-exponential-efficiency}
\leanverified{paper\_xi\_golden\_metallic\_perdigit\_exponential\_efficiency}
沿用定义 \ref{def:const-type} 与定理 \ref{thm:metallic-gap} 的记号，并把
\[
h(A):=\frac{\log\lambda_A}{A}
\qquad (A\in[1,\infty))
\]
视为实变量 \(A\) 上的解析函数。则：
\begin{enumerate}
  \item \(h(A)\) 在 \([1,\infty)\) 上严格递减；
  \item 等价地，\(\lambda_A^{1/A}\) 在 \([1,\infty)\) 上严格递减；
  \item 对每个 \(A>1\) 都有严格不等式
  \[
  \lambda_A<\varphi^A,
  \]
  特别地，对所有整数 \(A\ge 2\) 都成立。
\end{enumerate}
因此在全部常数型金属平均口径中，黄金端点 \(A=1\) 唯一最大化“每个连分位数所产生的指数增长效率”。
\end{corollary}
\begin{proof}
由定理 \ref{thm:metallic-gap}，
\[
\kappa(A):=\frac{A}{\log\lambda_A}
\]
在 \([1,\infty)\) 上严格递增且恒为正。故其倒数
\[
h(A)=\kappa(A)^{-1}=\frac{\log\lambda_A}{A}
\]
严格递减，得到第一条。

对正函数 \(h(A)\) 逐点取指数，便得
\[
\lambda_A^{1/A}=e^{h(A)}
\]
在 \([1,\infty)\) 上亦严格递减，得到第二条。

最后，由 \(\lambda_1=\varphi\) 与第二条，对任意 \(A>1\) 都有
\[
\lambda_A^{1/A}<\lambda_1=\varphi.
\]
两边取 \(A\) 次幂即得
\[
\lambda_A<\varphi^A.
\]
证毕。
\end{proof}

\endinput
